\label{sec:background}

\subsection{The Pre-\textit{Shelby} Section~5 Framework}

Section~5 of the Voting Rights Act required ``covered'' jurisdictions---
those designated under Section~4(b)'s coverage formula---to obtain
federal approval (``preclearance'') before implementing any change to
a voting law or procedure. The coverage formula in Section~4(b) was
based on two criteria: (1) use of a ``test or device'' (such as a
literacy test) as a condition of voting as of November 1964, 1968,
or 1972, and (2) voter registration or turnout below 50\% in the
presidential elections of those years.

The nine states fully covered under Section~5's original formula
were Alabama, Alaska, Arizona, Georgia, Louisiana, Mississippi,
South Carolina, Texas, and Virginia. Parts of additional states
were also covered, including counties in California, Florida,
Michigan, New Hampshire, New York, North Carolina, and South Dakota.

Preclearance required covered jurisdictions to demonstrate to
the Department of Justice (administrative preclearance) or the
U.S.\ District Court for the District of Columbia (judicial preclearance)
that any proposed change in voting procedure had neither the purpose
nor the effect of retrogressing minority voting strength. The
``non-retrogression'' standard, articulated in
\textit{Beer v.\ United States},\footnote{\textit{Beer v.\ United States},
425 U.S.\ 130 (1976).} meant that covered states could change their
voting procedures and redistrict, but they could not reduce the
minority electoral opportunity they had previously afforded.

For redistricting, Section~5 operated as an ex ante floor: each
redistricting cycle, covered states had to demonstrate that the new
plan did not reduce the number of majority-minority districts or
otherwise retrogress minority voting strength. This generated extensive
DOJ review of redistricting plans, with conditions and objections
imposed when plans fell below the non-retrogression standard.

\subsection{\textit{Shelby County v.\ Holder} (2013): The Holding}

Chief Justice Roberts, writing for the five-Justice majority, held
that Section~4(b)'s coverage formula was unconstitutional under the
Fifteenth Amendment's enforcement clause and the equal sovereignty
principle of the states.\footnote{\textit{Shelby County v.\ Holder},
570 U.S.\ 529, 544 (2013) (``The Fifteenth Amendment is not designed
to punish for the past; its purpose is to ensure a better future.'').}
The majority's reasoning rested on two premises.

First, the coverage formula was based on data from 1964--1972 and had
not been updated in the 2006 reauthorization, which simply extended the
existing formula for twenty-five years. The majority held that this was
constitutionally impermissible: the formula must be ``sufficiently
related to the problem it targets'' and the current problem of minority
vote suppression in formerly covered states was, in the majority's view,
not demonstrably worse than in non-covered states using 2012 data.\footnote{
\textit{Id.}\ at 556.}

Second, the majority applied the equal sovereignty principle---the
doctrine that states must be treated equally by federal legislation---and
concluded that Section~5's differential treatment of covered states
could not be sustained without a coverage formula calibrated to
current conditions.\footnote{\textit{Id.}\ at 554--56.}

The dissent by Justice Ginsburg (joined by Justices Breyer, Sotomayor,
and Kagan) argued that the majority had evaluated the VRA against its
own success: the coverage formula's jurisdictions had fewer Section~2
violations not because discrimination was absent but because Section~5
was deterring it.\footnote{\textit{Id.}\ at 591--92 (Ginsburg, J., dissenting).}

\subsection{Section~5 After \textit{Shelby County}: A Dead Letter}

Section~5 remains in the U.S.\ Code at 52~U.S.C.~\S~10304. But without
a valid coverage formula in Section~4(b), no jurisdiction can be
``covered'' under Section~5. The Supreme Court left Section~5 intact
as statutory text while making it unenforceable, in the expectation
that Congress would enact a new coverage formula.

Congress has not done so. The John Lewis Voting Rights Advancement Act,
which would update the coverage formula based on a new methodology
tracking voting rights violations in the preceding twenty-five years,
has passed the House of Representatives but has not cleared the Senate
as of 2026. The most powerful provision of the VRA---the one that
operated ex ante, with an administrative process that was faster and
cheaper than federal litigation---thus remains unenforceable.

\subsection{Immediate Post-\textit{Shelby} Redistricting Consequences}

The immediate redistricting consequence of \textit{Shelby County} was
the implementation of plans that had been held under preclearance review
or that would have been blocked under preclearance. Texas implemented
an interim congressional plan drawn by a three-judge court during
Section~5 litigation; after \textit{Shelby County}, the state enacted
its own plan. North Carolina redrew its congressional map in ways that
federal courts later found in \textit{Covington v.\ North Carolina}
to involve racial predominance in violation of the Equal Protection
Clause.\footnote{\textit{Covington v.\ North Carolina}, 316 F.R.D. 117
(M.D.N.C.\ 2016) (finding that 28 of 170 state legislative districts were
racially gerrymandered); see also \textit{North Carolina v.\ Covington},
581 U.S.\ 1 (2017) (per curiam).}

These consequences illustrate the prophylactic function that Section~5
had served: it did not merely remedy violations after the fact but
deterred them before enactment. Once that deterrent was removed, some
former covered states moved quickly to implement maps that would face
years of Section~2 litigation.
